\begin{figure}[H]
\centering
\begin{tikzpicture}
\begin{axis}[
    ybar stacked,
    bar width=36pt,
    width=14cm,
    height=9cm,
    ymin=0, ymax=11,
    ytick={0,2,4,6,8,10,11},
    ylabel={Number of responses},
    xlabel={Animation},
    xtick=data,
    symbolic x coords={A1, A2, A3, A4, A5, A6},
    xticklabels={
        {Curiosity},
        {Happiness},
        {Sadness},
        {Frustrated},
        {Anger},
        {Fear}
    },
    x tick label style={
        font=\small,
        rotate=20,
        anchor=east,
        yshift=-5pt,
    },
    legend style={
        at={(0.5,-0.28)},
        anchor=north,
        legend columns=3,
        font=\small,
        column sep=10pt,
    },
    enlarge x limits=0.12,
    tick label style={font=\small},
    label style={font=\small},
    xlabel style={yshift=-10pt},
    ymajorgrids=true,
    grid style={line width=0.2pt, draw=gray!30},
]

\addplot+[ybar, fill=intended, draw=white, line width=0.5pt]
    coordinates {
        (A1, 4)(A2, 11)(A3, 10)(A4, 9)(A5, 11)(A6, 8)
    };

\addplot+[ybar, fill=adjacent, draw=white, line width=0.5pt]
    coordinates {
        (A1, 5)(A2, 0)(A3, 1)(A4, 2)(A5, 0)(A6, 2)
    };

\addplot+[ybar, fill=other, draw=white, line width=0.5pt]
    coordinates {
        (A1, 2)(A2, 0)(A3, 0)(A4, 0)(A5, 0)(A6, 1)
    };

\legend{
    Intended emotion,
    Adjacent/related emotion,
    Other
}

\end{axis}
\end{tikzpicture}
\caption{Distribution of emotion recognition responses across all six
animations for the live study (N=11). Responses were grouped
into three categories: intended emotion (matching the target),
adjacent/related emotion (overlapping emotional state), and other
(unrelated responses). Dark purple indicates correct identification
of the intended emotion.}
\label{fig:emotion_recognition_live}
\end{figure}